\sect{客观题答案}
\begin{center}
\begin{tabular}{c*{10}{c}}
题号 & 1 & 2 & 3 & 4 & 5 & 6 & 7 & 8 & 9 & 10\\
\hline
答案 & C & A & D & B & A & C & D & B & A & C\\[3pt]
题号 & 11 & 12 & 13 & 14 & 15 & 16 & 17 & 18 & 19 & 20\\
\hline
答案 & 对 & 错 & 对 & 错 & 对 & 对 & 错 & 错 & 对 & 错\\
\end{tabular}
\end{center}

\sect{一、单项选择题详解}
\begin{enumerate}[label=\arabic*.]
\item \textbf{C。} A 将当前验证折的信息用于估计预处理参数。B 同样利用了验证数据拟合预处理，且与训练数据独立估计会造成变换口径不一致。C 在每折隔离验证集，最终一次使用留出测试集，可以保持评估独立。D 用测试结果比较超参数，测试集承担了验证集的作用。
\item \textbf{A。} 准确率为 $(72+102)/200=0.87$，精确率为 $72/(72+8)=0.9$，召回率为 $72/(72+18)=0.8$，$F_1=144/(144+8+18)\approx0.8471$。B 互换了精确率和召回率。C 把两者的调和平均算成算术平均。D 把正类召回率 $0.8$ 误作整体准确率。
\item \textbf{D。} 阈值为 $0.70$ 时，两个正例中有一个判正，两个负例中也有一个判正，所以 $\mathrm{TPR}=\mathrm{FPR}=0.5$。阈值从高到低扫过四个分数，ROC 顶点依次为 $(0,0)$、$(0,0.5)$、$(0.5,0.5)$、$(0.5,1)$、$(1,1)$，面积为 $0.5\times0.5+0.5\times1=0.75$。A 把单个 ROC 点的坐标当作面积。B 漏计分数为 $0.80$ 的假阳性。C 把分数为 $0.50$ 的正例也计成真阳性。
\item \textbf{B。} 用 $H(P,Q)=-\sum_i P_i\log_2Q_i$ 和 $\KL(P\,\|\,Q)=H(P,Q)-H(P)$，四个未舍入结果依次为 $1.5,1.29879494,1.91951798,0.41951798$。A 用反向交叉熵 $H(Q,P)=1.75$ 并减去 $H(P)$，两个式子的方向不一致。C 从正向交叉熵扣除了 $H(Q)$。D 把整组的 $P,Q$ 对调了。
\item \textbf{A。} 记 $h(p)$ 为二元熵。信息增益为 $1-0.4h(1/4)-0.6h(2/3)\approx0.1245$，用分支熵 $h(0.4)\approx0.9710$ 归一后得增益比 $0.1282$。加权基尼为 $0.4(2\cdot 1/4\cdot 3/4)+0.6(2\cdot 2/3\cdot 1/3)=5/12$。B 把未归一的信息增益直接作为增益比。C 给出的是基尼减少量 $1/2-5/12$。D 对两个分支的基尼取了等权平均。
\item \textbf{C。} A 漏乘类别先验。B 的先验只在对数判别差中增加常数。C 的判别差二次项为 $-\bm x^\top(\Sigma_1^{-1}-\Sigma_2^{-1})\bm x/2$，相等时抵消，不等时保留。D 取对数使高斯的指数项变为显式二次项，不会消去它们。
\item \textbf{D。} 曼哈顿距离依次为 $3,4,7$，欧氏距离依次为 $3,\sqrt8,5$。A 把乙误作曼哈顿最近点。B 把甲误作欧氏最近点。C 对调了两种距离的判决。
\item \textbf{B。} Xavier 用 $2/(64+32)=1/48$，He 用 $2/64=1/32$。A 对调两种结果。C 漏掉两个公式分子中的因子 $2$。D 把 He 的扇入错换为扇出。
\item \textbf{A。} LeNet 的局部卷积与下采样适合手写字符识别，AlexNet 的主要做法包括使用 ReLU 和部分全连接层的 Dropout。B 把 ResNet、GoogLeNet 的设计加给了这两个网络。C 的图像块自注意力和门控循环都不是这两种网络的主要机制。D 忽略了两者的卷积结构，AlexNet 也不靠循环状态处理空间信息。
\item \textbf{C。} 原始形式每次更新 $\bm w\gets\bm w+y_i\tilde{\bm x}_i$，累加成题给对偶展开。在有界且以正间隔线性可分的条件下，Novikoff 定理给出 $(R/\gamma)^2$ 的错误更新次数上界。A 丢掉了定方向的标签。B 把可分条件下的保证推广到不可分数据。D 一个样本可在不同轮次反复触发更新，样本数不构成一般上界。
\end{enumerate}

\sect{二、判断题详解}
\begin{enumerate}[label=\arabic*.,start=11]
\item \textbf{对。} 固定 $x$，将预测以其训练集平均值为中心，将真实输出以条件均值为中心，展开平方。两个中心化变量均值为零且相互独立，交叉项期望为零，余下偏差平方、模型方差和条件噪声方差。
\item \textbf{错。} PCA 的目标是保留方差，使用协方差最大的方向，与标签无关。LDA 则利用类标衡量类间分离与类内散布，可能保留低总方差但判别力较强的方向。
\item \textbf{对。} 早停用验证结果选轮次，增强扩展有效训练样本，单独的权重衰减可写成权重乘 $1-\eta\lambda$，其中 $0<\eta\lambda<1$。
\item \textbf{错。} 更新为 $w_{\mathrm{new}}=(1-\eta)w$。$\eta=0.01$ 时每步乘 $0.99$，减小较慢。$\eta=2.2$ 时每步乘 $-1.2$，符号来回反转且绝对值增大。
\item \textbf{对。} 连乘分别是 $0.5^L$ 与 $1.5^L$。范数裁剪只在梯度范数超过上限时缩小它，不能恢复因链式连乘而变得很小的梯度。
\item \textbf{对。} 堆叠的循环层同时有沿时间的状态输入和来自下层的当前时刻输入。三项任务分别由一个条件输入生成多步输出、由多步输入产生一个分类结果、为每步输入产生一个标签。
\item \textbf{错。} 题设使重构期望与近似后验无关，最大化 ELBO 此时只需让近似后验和先验的 KL 最小。取近似后验等于先验即可，隐变量可能丢失关于该样本的信息。
\item \textbf{错。} 虽然损失写成 $-\log p_y$，但 $p_y=\exp(z_y)/\sum_j\exp(z_j)$ 的分母含所有 logits。提高一项错误类 logit 会降低 $p_y$ 并增大损失。
\item \textbf{对。} 严格递增变换保持分数次序和并列关系。扫描所有阈值时产生相同的判正集合，ROC 及其面积随之保持。
\item \textbf{错。} 密度的约束是非负且积分为 $1$，密度值可以大于 $1$。例如标准差为 $0.1$ 的高斯密度在均值处为 $1/(0.1\sqrt{2\pi})>1$。
\end{enumerate}

\sect{三、计算与分析题参考答案与评分标准}

\noindent 等价的公式和解法给同等分数。数值合理舍入不扣分。前问数值错误且后问沿用时，后问按其步骤给分。每问扣分不超过该问分值。

\noindent\textbf{1. 多层感知机与反向传播}\pts{10}

\begin{enumerate}[label=(\arabic*)]
\item $\bm{z}=(r,-r)^{\mathsf T}$，由 $\tanh r=(e^{2r}-1)/(e^{2r}+1)$，得 $\bm{h}=(1/2,-1/2)^{\mathsf T}$。
线性部分和激活后向量各 $0.5$ 分。将 $\tanh$ 当成 Sigmoid，扣激活部分 $0.5$ 分。
\item $a=2(1/2)-(-1/2)+c=\ln 3$，所以 $p=3/4$，$L=-\ln p=\ln(4/3)\approx0.2877$。
得分、概率计算合计 $0.5$ 分，损失 $0.5$ 分。漏掉输出偏置，扣第一部分 $0.5$ 分。
\item
\[
\frac{\partial L}{\partial a}
=\left(-\frac{y}{p}+\frac{1-y}{1-p}\right)p(1-p)=p-y=-\tfrac14.
\]
写出两个因子并按链式法则相乘得 $0.5$ 分，化简与数值合计 $0.5$ 分。符号写成 $y-p$，扣后者 $0.5$ 分。
\item $\nabla_{\bm{v}}L=(p-y)\bm{h}=(-1/8,1/8)^{\mathsf T}$，$\partial L/\partial c=p-y=-1/4$。
两项各 $0.5$ 分。将输入 $\bm{x}$ 代替 $\bm{h}$，扣权重梯度 $0.5$ 分。
\item 隐藏层收到的梯度为 $\partial L/\partial\bm{h}=(p-y)\bm{v}$。逐元素乘上 $1-h_i^2=3/4$，得
\[
\frac{\partial L}{\partial\bm{z}}=(-3/8,3/16)^{\mathsf T}.
\]
正确传回输出权重得 $0.5$ 分，乘激活导数并算对得 $0.5$ 分。漏掉 $\tanh$ 的导数，扣后者 $0.5$ 分。
\item 令 $\bm{\delta}=\partial L/\partial\bm{z}$，则
\begin{align*}
\nabla_{\bm{W}}L=\bm{\delta}\bm{x}^{\mathsf T}
&=\tfrac1{16}\begin{pmatrix}-6&6&-12\\3&-3&6\end{pmatrix},\\
\nabla_{\bm{b}}L&=\tfrac1{16}(-6,3)^{\mathsf T}.
\end{align*}
矩阵和偏置各 $0.5$ 分。将外积写反或矩阵维度不符，扣矩阵部分 $0.5$ 分。
\item 对 $g_c=-1/4$，未修正的一阶、二阶矩为 $m_c=(1-\beta_1)g_c=-1/8$、$s_c=(1-\beta_2)g_c^2=1/64$。偏差修正后 $\widehat m_c=m_c/(1-\beta_1)=-1/4$，$\widehat s_c=s_c/(1-\beta_2)=1/16$。
未修正和修正各 $0.5$ 分。用平方的一阶矩替代二阶矩或漏修正，扣对应 $0.5$ 分。
\item 第一步的修正矩对任意参数均为 $g,g^2$。设某坐标的原梯度为 $k/16$，则该坐标改变量为 $-k/[20(|k|+1)]$。故
\begin{align*}
\bm{W}'&=\bm{W}+\tfrac1{20}\begin{pmatrix}6/7&-6/7&12/13\\-3/4&3/4&-6/7\end{pmatrix},\\
\bm{b}'&=\tfrac1{20}(6/7,-3/4)^{\mathsf T}.
\end{align*}
两组参数各 $0.5$ 分。把 $\epsilon$ 放进根号会得到错误更新量，扣涉及的对应部分。
\item $\bm{v}'=(2,-1)^{\mathsf T}+\tfrac1{20}(2/3,-2/3)^{\mathsf T}$，$c'=c+\tfrac1{20}(4/5)=\ln 3-\tfrac32+\tfrac1{25}$。
两组参数各 $0.5$ 分。将更新写成沿正梯度移动，扣方向写错的对应部分。
\item 四组梯度的平方和合计 $[270+45+8+16]/16^2=339/256$。因此缩放因子为 $16/\sqrt{339}\approx0.8690$，裁剪后 $g_c=-4/\sqrt{339}\approx-0.2173$。
范数及统一因子 $0.5$ 分，裁剪后的值 $0.5$ 分。逐元素截断或将平方和直接当范数，扣对应部分。
\end{enumerate}

\noindent\textbf{2. 支持向量机与主成分特征}\pts{8}
\begin{enumerate}[label=(\arabic*)]
\item $\bar{\bm x}=(0,0)^\top$（0.5 分）。四项外积求和再除以 4，
$S=\begin{pmatrix}6&4\\4&6\end{pmatrix}$（0.5 分）。均值错误扣 0.5 分，使用 $1/3$ 或协方差任一项错误扣矩阵的 0.5 分。
\item 特征值为 $10,2$，第一主轴为 $\bm u=(1,1)^\top/\sqrt{2}$（0.5 分）。四个投影坐标（0.5 分）为
\[
(z_1,z_2,z_3,z_4)=(-3\sqrt{2},-\sqrt{2},\sqrt{2},3\sqrt{2}).
\]
取小特征值对应方向或遗漏单位化扣前 0.5 分，错算投影扣后 0.5 分。
\item 投影方差为 $\bm u^\top S\bm u$（0.5 分）。令
$\mathcal J=\bm u^\top S\bm u-\lambda(\bm u^\top\bm u-1)$，对 $\bm u$ 求导得 $S\bm u=\lambda\bm u$（0.25 分）。乘以 $\bm u^\top$ 得投影方差等于对应特征值，故取最大特征值（0.25 分）。漏平方、漏单位约束、选错特征值，分别扣对应分数。
\item 约束为 $y_i(wz_i+b)\ge1$（0.5 分）。由内侧两点得到 $\sqrt{2}w\pm b\ge1$，所以 $w\ge1/\sqrt{2}$。取 $w=1/\sqrt{2}, b=0$ 可行且使 $w^2/2$ 最小，$f(z)=z/\sqrt{2}$，总宽度 $2/|w|=2\sqrt{2}$（0.5 分）。漏掉间隔约束扣前 0.5 分，函数或宽度错误扣后 0.5 分。
\item $w=\sum_i\alpha_i y_i z_i$，$\sum_i\alpha_i y_i=0$（各 0.25 分）。条件里漏掉 $y_i$ 或遗漏偏置的平稳条件，扣对应的 0.25 分。
\item 两个外侧点的函数间隔为 3，互补松弛给 $\alpha_1=\alpha_4=0$（0.5 分）。等式约束给 $\alpha_2=\alpha_3$，由 $w=2\sqrt{2}\alpha_2$ 得 $(0,1/4,1/4,0)$，所以支持向量是第 2、3 个点（0.5 分）。给外侧点非零乘子，或内侧乘子及支持向量判断错误，扣对应的 0.5 分。
\item $z=-3/\sqrt{2}$，$f(z)=-3/2$，判为 $-1$。漏投影或实际判别值错误扣 0.25 分，类别错误扣 0.25 分。
\item 各乘子均在 $[0,1/8]$ 内，且两非零乘子的标签相反，$\sum_i\alpha_iy_i=0$（0.5 分）。$w=2\sqrt{2}/8=\sqrt{2}/4$（0.5 分）。漏检任一约束扣 0.5 分，把 $C$ 直接当 $w$ 或尺度错误扣 0.5 分。
\item 四点函数间隔依次为 $(1.5,0.5,0.5,1.5)$。松弛变量为
\[
\xi_i=\max\{0,1-y_if(z_i)\},\qquad
(\xi_1,\xi_2,\xi_3,\xi_4)=(0,0.5,0.5,0).
\]
第 2、3 个点在间隔内，分类仍正确。松弛变量错误扣 0.25 分，把位于间隔内判成错分或选错点扣 0.25 分。
\item $\varphi(z)=z^2$，有 $z^2(z')^2=(zz')^2$。映射或乘积关系错误扣 0.5 分。
\end{enumerate}

\noindent\textbf{3. 二维测点的密度估计与聚类}\pts{12}
\begin{enumerate}[label=(\arabic*)]
\item $\widehat p(q)=\frac15\sum_{u\in\{a,b,c,d,e\}}K_h(q-u)$。五个距离的平方为 $1,2,9,8,17$，所以
\[
\widehat p(q)=\frac{\ln 2}{5\pi}
 (2^{-1}+2^{-2}+2^{-9}+2^{-8}+2^{-17})\approx0.03335.
\]
\textbf{评分} 核平均式 0.5 分，代入及结果 0.5 分。漏 $1/5$ 或核的归一化系数，扣相应 0.5 分。
\item 曼哈顿距离依次为 $1,2,3,4,5$，前三个是 $a,b,c$，预测负类。欧氏距离平方依次为 $1,2,9,8,17$，前三个是 $a,b,d$，预测正类。
\textbf{评分} 每种距离的排序和投票合计 0.5 分。该种排序或投票错则扣 0.5 分。
\item 新估计为
\[
 \widehat p_{\mathrm{new}}(q)=\frac{5\widehat p(q)+3K_h(q-e)}8.
\]
原五点中，$e$ 离 $q$ 最远，所以 $K_h(q-e)$ 小于这五项的平均值 $\widehat p(q)$，新估计小于原估计。
\textbf{评分} 加权更新式 0.5 分，比较方向与依据 0.5 分。增加了三项却仍除以 5，扣更新式 0.5 分。
\item 未选点 $b,c,e$ 到较近中心的平方距离分别为 $1,5,5$，因此第三个中心落在 $\{b,e\}$ 的概率为 $(1+5)/(1+5+5)=6/11$。给定落在该集合，候选点只有 $b,e$，选到 $e$ 的概率为 $5/(1+5)=5/6$。
\textbf{评分} 两个概率各 0.5 分。只计算到单个中心的距离或给定事件后未改分母，各扣受影响问的 0.5 分。
\item $a,b$ 分到以 $a$ 为中心的簇，$c,d,e$ 分到以 $d$ 为中心的簇。平方误差和为 $0+1+5+0+5=11$。
\textbf{评分} 分配 0.5 分，误差和 0.5 分。将距离而非距离平方相加，扣 0.5 分。
\item 新中心为 $(1,1/2)$ 和 $(1,3)$。五个平方误差分别为 $1/4,1/4,1,2,1$，误差和为 $9/2$。
\textbf{评分} 两中心都对 0.5 分，误差和 0.5 分。仍用初始中心计算误差扣 0.5 分。
\item 以括号外的下标表示合并高度，单链树为 $(((U,V)_{4},W)_{6},Z)_{9}$。初次合并后，两条用于判定的最短跨簇距离为 $d(\{U,V\},W)=6$ 及 $d(\{U,V,W\},Z)=9$。
\textbf{评分} 完整合并次序 0.5 分，高度均正确 0.5 分。树图或等价的带高度嵌套表达均可。
\item 全链树为 $((U,V)_{4},(W,Z)_{9})_{19}$。因为合并 $U,V$ 后，到 $W$ 的全链距离是 10，大于 $d(W,Z)=9$。平均链树为 $(((U,V)_{4},W)_{8},Z)_{43/3}$。中间两步高度分别由 $(10+6)/2$ 和 $(19+15+9)/3$ 给出。
\textbf{评分} 每种方法的正确合并次序与高度各 0.5 分，共 2 分。把簇间平均误当成不同规模子簇的等权平均，扣该法高度项 0.5 分。
\item 五点邻域依次为 $\{a,b\},\{a,b,d\},\{c,e\},\{b,d\},\{c,e\}$。因此 $b$ 为核心点，$a,d$ 为边界点，$c,e$ 为噪声点。
\textbf{评分} 邻域及核心点 1 分，边界、噪声各 0.5 分。漏自身或漏距离等于阈值的边时，在邻域项扣 0.5 分，后续按所得邻域正确分类可给后续方法分。
\item 更新均值用加权平均 $\mu_k^{\mathrm{new}}=\sum_i r_{ik}x_i/\sum_i r_{ik}$。对非最近成分与最近成分的责任度比，指数中有严格为负的平方距离差除以 $2\sigma^2$，该比值趋于零，最近成分的责任度趋于 1。
\textbf{评分} 加权均值 0.5 分，极限及距离差的解释 0.5 分。反过来声称远点责任度趋于 1 扣 0.5 分。
\end{enumerate}

\noindent\textbf{4. 自编码器、证据下界与高斯变分后验}\pts{8}

\begin{enumerate}[label=(\arabic*)]
\item 确定性自编码器把给定图像映射到确定的编码向量。VAE 编码器给出近似后验的参数，例如均值和方差，再从该分布中抽取潜变量用于训练时的重构。生成新图像时从先验 $p(z)$ 抽样，再用解码器生成。
编码输出的区别与生成时使用先验各 $0.25$ 分。将生成时使用的先验写成依赖某个输入的近似后验，扣后 $0.25$ 分。
\item $p(x=1,0)=3/8$、$p(x=1,1)=1/8$，相加得证据 $1/2$。分别除以 $q(z\mid x=1)=1/2$，比值为 $3/4,1/4$。
边缘概率、两个比值各 $0.5$ 分。直接用两个条件概率之和算边缘概率，扣边缘概率部分。
\item 利用两点均有 $q(z\mid x)>0$，有
\[
p(x)=\sum_zq(z\mid x)\frac{p(x,z)}{q(z\mid x)}
=\E_q\!\left[\frac{p(x,z)}{q(z\mid x)}\right].
\]
对凹函数 $\ln$ 用 Jensen 不等式，得到
\[
\ln p(x)\ge\E_q\!\left[\ln\frac{p(x,z)}{q(z\mid x)}\right]
=\mathcal L(q).
\]
写对加权期望 $0.5$ 分，指出 $\ln$ 的凹性 $0.5$ 分，方向及下界表达式 $0.5$ 分。将对数的期望与期望的对数直接画等号，扣方向及下界部分。
\item 将对数内部拆开，得
\begin{align*}
\mathcal L(q)&=\E_q[\ln p(x\mid z)]+\E_q[\ln p(z)-\ln q(z\mid x)]\\
&=\E_q[\ln p(x\mid z)]-\KL(q(z\mid x)\Vert p(z)).
\end{align*}
重构项与带正确符号、方向的 KL 项各 $0.5$ 分。把 KL 写成相反顺序或加号，扣 KL 部分。
\item
\begin{align*}
\mathcal L(q)&=\tfrac12\ln(3/4)+\tfrac12\ln(1/4)
=\tfrac12\ln(3/16)\approx-0.8370,\\
\ln p(x=1)-\mathcal L(q)&=\ln(2/\sqrt3)\approx0.1438.
\end{align*}
下界、差值各 $0.5$ 分。因先验 KL 为零就把下界当成对数证据，扣差值部分。
\item 由于 $\ln$ 严格凹，在 $q$ 支持两个状态时，等号要求比值 $p(x,z)/q(z\mid x)$ 在两点相同。由归一化，该常数为 $p(x)$，所以 $q(z\mid x)=p(x,z)/p(x)$，两点依次为 $3/4,1/4$。
等号条件和具体结果各 $0.5$ 分。写成 $q=p(z)$，扣具体结果部分。
\item
\[
\ln\frac{q_c(z)}{p_c(z)}=-\tfrac12\ln v-\frac{(z-\mu)^2}{2v}+\tfrac12z^2.
\]
三项符号及系数正确得 $0.5$ 分。将方差 $v$ 当成标准差代入，扣 $0.5$ 分。
\item $\E_{q_c}[(z-\mu)^2]=v$，$\E_{q_c}[z^2]=v+\mu^2$，代入得
\[
\KL(q_c\Vert p_c)=\tfrac12(v+\mu^2-1-\ln v).
\]
二阶矩 $0.5$ 分，闭式结果 $0.5$ 分。漏掉均值平方，扣受影响部分。
\item 只在 $z=1$ 取期望时，期望值为 $p(x=1,z=1)=1/8$，漏掉 $z=0$ 对证据的贡献 $3/8$，所以两者不等。指出实际期望及漏项得 $0.5$ 分。把零概率点也代入比值计算，不给该问分。
\end{enumerate}

\noindent\textbf{5. 卷积、逆向序列网络与掩码注意力}\pts{8}

\begin{enumerate}[label=(\arabic*)]
\item 无空洞时，输入一维大小为 $N$、两侧各填充 $p$、核宽 $k$、步长 $s$，输出为 $\lfloor(N+2p-k)/s\rfloor+1$。两方向分别代入 $N=5,6$，得到 $3\times4\times3$。
公式 $0.5$ 分，三维尺寸 $0.5$ 分。最后一维沿用输入的 $2$，扣尺寸部分。
\item 两通道的加权和分别为 $1+0+0+1=2$ 和 $0+2-1+0=1$。加上一次偏置得到 $7/2$。可训练参数量为 $3(2\cdot2\cdot2+1)=27$。
输出与参数量各 $0.5$ 分。按输入通道重复加偏置，扣输出部分。不计输出通道偏置，扣参数量部分。
\item 空间尺寸仍为 $3\times4$，输出为 $3\times4\times4$，参数为 $4(3+1)=16$。尺寸和参数量各 $0.25$ 分。把卷积核深度当作 $1$，扣参数量部分。
\item 雅可比为 $I+J_F(\bm{x})$，反传中 $I^{\mathsf T}$ 提供了一条不经过 $F$ 的梯度路径。令 $F$ 为零可让新增块实现恒等映射，因而能保留较浅模型的映射。这些机制有助于训练，梯度的总和仍由两条路径共同决定。
雅可比和梯度路径 $0.25$ 分，恒等映射对退化的解释 $0.25$ 分。只写“网络更深所以梯度更大”，不给分。
\item $o_t=\sigma(\bm{v}^{\mathsf T}\bm{h}_t)$，其中 $\bm{h}_t$ 由当前输入和 $\bm{h}_{t+1}$ 算出。先得 $\bm{h}_3=(2,-1)^{\mathsf T}$，再得 $\bm{h}_2=(2,-1/3)^{\mathsf T}$，故 $o_2=\sigma(2)=1/(1+e^{-2})\approx0.8808$。
明确从右向左及输出形式 $0.5$ 分，状态与输出数值 $0.5$ 分。误用 $\bm{h}_{t-1}$，扣递推与相应数值部分。
\item 记局部梯度为 $\bm{g}_t$、汇总梯度为 $\bm{\delta}_t$。由于 $\bm{h}_t$ 会输入 $\bm{h}_{t-1}$，梯度按 $t=1,2,3$ 递推，$\bm{\delta}_0=\bm{0}$ 且
\[
\bm{\delta}_t=\bm{g}_t+\bm{W}^{\mathsf T}\bm{\delta}_{t-1}.
\]
得到 $\bm{\delta}_1=(1,0)^{\mathsf T}$、$\bm{\delta}_2=(1/2,1)^{\mathsf T}$、$\bm{\delta}_3=(1/4,1/3)^{\mathsf T}$。同一个 $W_{11}$ 在三个时刻使用，梯度为 $\sum_{t=1}^3\delta_{t,1}h_{t+1,1}=1\cdot2+(1/2)\cdot2+(1/4)\cdot0=3$。
递推方向和公式 $0.5$ 分，参数梯度的累加与数值 $0.5$ 分。仅使用一个时刻的梯度，扣参数梯度部分。
\item 一步直接导数为该维的 $f_t$，跨 $k$ 步得到 $\prod_{j=t-k+1}^t f_j$。若相关的遗忘门值能长时间接近 $1$，沿此路径的衰减就较慢。
梯度乘积与解释各 $0.5$ 分。乘上输出门或把遗忘门相加，扣梯度部分。
\item 采用行向量，第 $j$ 头的 $W_j^Q,W_j^K,W_j^V$ 都为 $8\times4$，每头输出 $n\times4$。拼接为 $n\times8$，再乘 $8\times8$ 的输出投影。总量为 $2\cdot3\cdot8\cdot4+8\cdot8=256$。
尺寸 $0.5$ 分，参数量 $0.5$ 分。漏掉输出投影或把两个头先求均值，扣对应部分。
\item 未来位置的得分加 $-\infty$。对尚可见的指数值逐行归一化，得权重矩阵
\[
\begin{pmatrix}1&0&0\\2/3&1/3&0\\1/6&1/2&1/3\end{pmatrix}.
\]
掩码方向及归一化均正确得 $0.5$ 分。将被掩码的得分置零后仍对其取指数，扣该问 $0.5$ 分。
\item 有 $n^2$ 对查询与键。每对点积的代价按头维度计，再对头求和为 $O(n^2d)$。对值向量加权求和同阶。点积、加权求和各 $0.25$ 分。只计算单个位置的代价，扣对应部分。
\end{enumerate}

\noindent\textbf{6. 工单分类中的朴素贝叶斯}\pts{8}
\begin{enumerate}[label=(\arabic*)]
\item $P(\text{急}\mid F)=4/7,\ P(\text{室外}\mid F)=4/6$。$P(\text{急}\mid N)=2/7,\ P(\text{室外}\mid N)=2/6$。两个分母分别使用 $4+3$ 和 $4+2$。
\textbf{评分} 每类的两个概率各合计 0.5 分。将两个特征都除以 6 或 7，扣受影响类的分。
\item 条件似然为 $8/21$、$2/21$。先验相同，故 $P(F\mid\text{急、室外})=(8/21)/(10/21)=4/5$，判为 $F$。
\textbf{评分} 对应类条件概率相乘 0.5 分，归一化与分类 0.5 分。把似然直接当后验扣 0.5 分。
\item $F$ 的词概率向量是 $(6,3,2)/11$，$N$ 是 $(2,3,6)/11$。两类各出现了 8 个词元，词表大小是 3，所以分母为 $8+3$，分子是该词在对应类别中的累计次数加 1。
\textbf{评分} 概率向量 0.5 分，平滑分母解释 0.5 分。用文档数作分母扣 0.5 分。
\item 两个“振动”分别贡献 $\ln(6/2)$，“发热”的条件概率在两类相同。相同先验相消，后验对数优势为 $2\ln 3$。赔率为 9，故 $P(F\mid x)=9/10$。
\textbf{评分} 对数优势 0.5 分，后验概率 0.5 分。重复词只计一次，扣对数优势项 0.5 分，后验按自己的赔率正确转换仍可给方法分。
\item 文档的存在指示为 $(1,1,0)$。$F$ 下似然为 $(4/6)(3/6)(1-2/6)=2/9$，$N$ 下为 $(2/6)(3/6)(1-4/6)=1/18$，故判为 $F$。
\textbf{评分} 两类非零特征及缺席因子都正确 0.5 分，似然比较与分类 0.5 分。漏乘“异响”未出现的因子扣 0.5 分。
\item 公式为 $\widehat\mu=n^{-1}\sum_i x_i,\ \widehat\sigma^2=n^{-1}\sum_i(x_i-\widehat\mu)^2$。$F$ 的均值为 $(3,2)$，方差为 $(1,1)$。$N$ 的均值为 $(1,1)$，方差为 $(1,4)$。
\textbf{评分} 两类均值合计 0.5 分，方差合计 0.5 分。用无偏方差的 $n-1$ 作分母扣方差项 0.5 分。
\item 两个特征的标准化平方残差在 $F$ 中相加为 $1$，在 $N$ 中相加为 $1+1/4$。标准差乘积之比为 $1:2$，相同先验相消，故对数优势 $\ln 2+1/8$。后验概率为 $(1+\tfrac12 e^{-1/8})^{-1}\approx0.6938$，判 $F$。
\textbf{评分} 正确的对数优势 0.5 分，后验与分类 0.5 分。忽略两类方差不同导致的正态化因子，扣对数优势项 0.5 分。
\item 朴素贝叶斯拟合 $P(Y)$ 和各类下的特征分布 $P(X_j\mid Y)$，再用贝叶斯公式得到后验。逻辑回归直接拟合 $P(Y\mid X)$。对正因子取对数会把许多很小的概率相乘改成相加，避免相乘下溢，同时保持得分大小次序。
\textbf{评分} 两种建模对象 0.5 分，对数作用 0.5 分。将两种建模对象颠倒扣 0.5 分。
\end{enumerate}

\noindent\textbf{7. 传感器校准、回归与正则化}\pts{8}
\begin{enumerate}[label=(\arabic*)]
\item $X=\begin{pmatrix}1&-1\\1&0\\1&1\end{pmatrix}$，$\bm y=(1,1,3)^\top$。漏截距列或矩阵错误扣 0.25 分，目标向量错误扣 0.25 分。
\item $X^\top X\bm w=X^\top\bm y$（0.5 分）。$X^\top X=\begin{pmatrix}3&0\\0&2\end{pmatrix}$，$X^\top\bm y=(5,2)^\top$（各 0.25 分）。漏写方程或错算矩阵及向量，扣对应分数。
\item 解得 $w_0=5/3,\ w_1=1$（各 0.5 分），预测函数 $\hat y=5/3+x$。某一系数与预测函数中对应项错误，扣该项的 0.5 分。
\item $x=0$ 给出 $v_0=1$，另外两点给出 $-v_1+v_2=0$ 与 $v_1+v_2=2$（0.5 分）。得 $v_1=v_2=1$，即 $\hat y=1+x+x^2$（0.5 分）。未用三点建立正确方程扣前 0.5 分，解错扣后 0.5 分。
\item 线性模型在两点的预测是 $(-1/3,11/3)$，残差是 $(-4/3,-1/3)$。$\mathrm{MSE}=17/18\approx0.9444$，$\mathrm{MAE}=5/6\approx0.8333$。
二次模型预测 $(3,7)$，残差 $(2,3)$。$\mathrm{MSE}=13/2=6.5$，$\mathrm{MAE}=5/2=2.5$。四个指标各 0.25 分。漏取绝对值、漏平方或误用分母，扣相应指标的分数。
\item 线性模型的两项验证误差都更低，二次模型虽然在三点上训练误差为 0，却在验证点的误差更大。选错模型扣 0.25 分，未说明训练与验证的差别扣 0.25 分。
\item 依赖参数的项为 $\frac{1}{2\sigma^2}\sum_i(y_i-w_0-w_1x_i)^2$（0.5 分）。前面的系数是固定正数，最小化该负对数似然与最小化残差平方和有相同解（0.5 分）。系数或平方残差写错扣前 0.5 分，未说明等价理由扣后 0.5 分。
\item $\sigma(\ln 2)=1/(1+1/2)=2/3<0.75$，判为不报警。概率错误扣 0.25 分，未用指定阈值或分类错误扣 0.25 分。
\item $L(t)=\sum_i\ln(1+e^{-tm_i})$（0.5 分）。每个 $m_i>0$，故 $t\to\infty$ 时损失趋于 0，但每个有限 $t$ 的损失均大于 0。参数范数沿该射线无界（0.5 分）。错用指数符号扣前 0.5 分，把下确界当有限参数处的最小值扣后 0.5 分。
\item 令 $\bm\theta=(\bm w,b)$，惩罚项为某正系数乘 $t^2\|\bm\theta\|^2$，随着 $t\to\infty$ 发散，总目标也发散。漏掉二次增长或误称放大可使正则化目标趋零，扣 0.5 分。
\end{enumerate}

\noindent\textbf{8. 滤器状态的识别与更新}\pts{8}
\begin{enumerate}[label=(\arabic*)]
\item $\alpha_1(i)=\pi_i b_i(o_1)$，$\alpha_t(j)=b_j(o_t)\sum_i\alpha_{t-1}(i)a_{ij}$。$a_{ij}$ 为 $i$ 到 $j$ 的转移概率，$b_j(o_t)$ 为到达 $j$ 后发射当前观测的概率。
\textbf{评分} 初值、递推各 0.5 分。发射概率错乘在离开的状态上，扣递推分。
\item $\alpha_1=(1/4,1/8)$。
\textbf{评分} 分量各 0.25 分。漏乘发射概率扣该分量。
\item $\alpha_2=(7,5)/128$。例如第一分量为 $(\frac14\cdot\frac34+\frac18\cdot\frac14)\frac14$。
\textbf{评分} 分量各 0.25 分。漏乘发射概率扣该分量。
\item $\alpha_3=(13,22)/1024$，发射概率向量为 $(1/4,1/2)$。
\textbf{评分} 两个分量各 0.25 分。混用 $M$ 和 $H$ 的发射概率扣受影响分量的分。
\item $\alpha_4=(122,79)/16384$，发射概率向量为 $(1/2,1/4)$。
\textbf{评分} 两个分量各 0.25 分。错用早一刻的观测扣受影响分量的分。
\item $P(O\mid\lambda)=\sum_i\alpha_4(i)=201/16384\approx0.01227$。
\textbf{评分} 求和得 0.5 分，取最大扣全分。
\item $\beta_4=(1,1)$，$\beta_t(i)=\sum_j a_{ij}b_j(o_{t+1})\beta_{t+1}(j)$。故 $\beta_3=(7,5)/16$。
\textbf{评分} 终止条件和递推 0.25 分，$\beta_3$ 0.25 分。用 $o_t$ 代替 $o_{t+1}$，扣递推项 0.25 分。
\item $\beta_2=(31,37)/256$，$\beta_1=(65,71)/2048$。
\textbf{评分} 各 0.25 分，下标错扣该项。
\item $\gamma_2(i)=\alpha_2(i)\beta_2(i)/P(O\mid\lambda)$，从而 $\gamma_2=(217,185)/402$。
\textbf{评分} 公式、结果各 0.25 分。漏后向量扣公式分。
\item $\xi_t(i,j)=\alpha_t(i)a_{ij}b_j(o_{t+1})\beta_{t+1}(j)/P(O\mid\lambda)$。
$\xi_t(1,2)$ 依次为 $37/201$、$35/201$ 和 $13/201$。
\textbf{评分} 公式、数值各 0.5 分。数值错一处扣 0.25 分，扣完为止。边缘后验相乘扣公式分。
\item 对 $T=4$，三组更新式为
\[
 \pi_i'=\gamma_1(i),\qquad
 a_{ij}'=\frac{\sum_{t=1}^{T-1}\xi_t(i,j)}{\sum_{t=1}^{T-1}\gamma_t(i)},\qquad
 b_i'(v)=\frac{\sum_{t=1}^{T}\gamma_t(i)\mathbf{1}\{o_t=v\}}{\sum_{t=1}^{T}\gamma_t(i)}.
\]
前三步状态 1 的后验为 $130/201$、$217/402$、$91/201$，和为 $659/402$。于是 $a_{12}'=(85/201)/(659/402)=170/659$。
\textbf{评分} 各更新式及数值各 0.5 分。计入 $\gamma_4(1)$ 扣转移式的 0.5 分。
\end{enumerate}
